% !TEX root = ../print.tex
% Draft \S3 --- Axioms
%
% Source: draft/spatial-hemigroup-scale-space.md, \S3 (lines 77-138).
%
% The axiom system is the causal one with causality replaced by reflection symmetry; that
% is the whole difference between the two articles, and this chapter records it. The
% relation to the causal theory is stated once, in rem:provenance-causal, and nowhere else
% in this chapter.
%
% Two members are exhibited because the theory has to fit between them. The Gaussian clause
% is elementary and self-contained; the Matern clause needs its kernels to be symmetric
% probability measures, which is settled in Chapter 10, so the node carries forward \uses
% edges. Those edges are honest and the graph stays acyclic: Chapter 10 reaches this
% chapter only through def:cascade-family.
%
% prop:no-positivity-no-classification is the article's thesis in negative form and is the
% one node here with real work in it; its proof is the draft's, written out.

\chapter{Axioms}
\label{ch:axioms}

The primitive object is a two-parameter family of operators indexed by ordered pairs of
scales. Positivity is stated in its place in the list, but the theory is developed so that
its role stays visible: \ref{prop:no-positivity-no-classification} says what the other axioms
admit without it, and \ref{thm:main-characterization} what it cuts that down to.

% draft: Definition 3.1
% twin: hcs def:cascade-family (Hemigroup.CascadeFamily) --- identical except that (A3)
% reads reflection symmetry here and causal support there, and that the ambient space is
% L^1(R) rather than L^1 of a half-line.
% The Lean rendering splits the axiom list: PreCascadeCore carries (A1)--(A3) and (A5)--(A7),
% IsPositive is (A4) and IsNondegenerate is (ND), with CascadeCore extending PreCascadeCore by
% the two. The split is what makes prop:no-positivity-no-classification -- a family satisfying
% every axiom but (A4) -- statable at all, and it lets each later node carry exactly the axiom
% list this definition gives it. Nothing about the definition itself changes.
\begin{definition}[symmetric cascade measurement family]\label{def:cascade-family}
\lean{SpatialLine.PreCascadeCore, SpatialLine.IsPositive, SpatialLine.IsNondegenerate,
SpatialLine.CascadeCore, SpatialLine.IsScaleCovariant, SpatialLine.CascadeFamily}\leanok
\notes{scale-space-axioms}
A family of linear operators $\Phis{s}{t} : \Lone \to \Lone$, $0 \le s \le t < \infty$, is a
\emph{symmetric cascade measurement family} if:
\begin{enumerate}
\item[(A1)] \emph{Boundedness.} Each $\Phis{s}{t}$ is a bounded operator on $\Lone$.
\item[(A2)] \emph{Translation covariance.} $\Phis{s}{t}\,\Tr{a} = \Tr{a}\,\Phis{s}{t}$ for
every $a \in \RR$.
\item[(A3)] \emph{Reflection symmetry.} $\Phis{s}{t}\,R = R\,\Phis{s}{t}$.
\item[(A4)] \emph{Positivity.} $\Phis{s}{t}\,L^1_+ \subseteq L^1_+$.
\item[(A5)] \emph{Unit mass.} $\int \Phis{s}{t} f = \int f$ for every $f \in L^1_+$.
\item[(A6)] \emph{Cascade (hemigroup).} $\Phis{t}{t} = \mathrm{Id}$, and
$\Phis{s}{t}\,\Phis{r}{s} = \Phis{r}{t}$ for all $r \le s \le t$.
\item[(A7)] \emph{Continuity.} For every $f \in \Lone$ the map $(s,t) \mapsto \Phis{s}{t}f$ is
continuous from $\{0 \le s \le t\}$ to $\Lone$.
\item[(A8)] \emph{Scale covariance.} There is a family $(S_\lambda)_{\lambda > 0}$ of
increasing bijections of $[0,\infty)$ with
$\Dil{\lambda}\,\Phis{s}{t} = \Phi_{S_\lambda s,\,S_\lambda t}\,\Dil{\lambda}$ for all
$0 \le s \le t$ and all $\lambda > 0$.
\item[(ND)] \emph{Nondegeneracy.} $\Phis{s}{t} \ne \mathrm{Id}$ whenever $s < t$.
\end{enumerate}
\statusT
\end{definition}

\begin{remark}[provenance: Pauwels et al.]\label{rem:provenance-pauwels}
The axioms of \sscite{pauwels1995extended} are, for a family $\{K_t\}_{t \ge 0}$: a linear
integral operator with an $\Lone$ kernel, shift invariance, evenness of the kernel, unit mass,
continuity of the kernel in $x$ and in $t$, \emph{recursivity} $K_sK_t = K_{s+t}$, and scale
invariance with a free rescaling function, $k_t(x) = \psi(t)^{-1}\phi(x/\psi(t))$ with $\psi$
continuous, strictly increasing, $\psi(0) = 0$ and $\psi(t) \to \infty$. Their linearity,
integral form and shift invariance are (A1)--(A2) together with the representation of
Chapter~\ref{ch:representation}. Evenness is (A3), unit mass is (A5), and continuity in $t$ is
(A7). Recursivity is (A6) with the tacit assumption that the increments are stationary in the
scale parameter, which is what this article drops. Their $\psi$ is the relabelling $S_\lambda$
of (A8), recovered rather than posited in \ref{prop:canonical-gauge}. Two things they assume
and this article does not: kernels that are continuous \emph{functions}, and the semigroup.
Measures are allowed here so that the identity is the diagonal limit, and nothing is lost,
since by \ref{prop:kernel-regularity} every nondegenerate admissible kernel has a density
anyway. Positivity, which in their paper enters only at the end, to bound the stability index
through a second moment, is an axiom here.
\end{remark}

% CHANGED (article mirror 2026-09-14, R157): three wordings, each in the safe direction. "The two
% systems differ in nothing else" is dropped -- one axiom is weakened and two smaller changes go
% with it, which is what the remark itself lists. "The degenerate member is not a pure translation"
% is replaced by "the member that stands where the pure delay stood": the identity family is
% excluded by (ND) and not by (A3), and calling the Gaussian the degenerate member contradicts the
% word "nondegenerate" used of it everywhere else. And the causal article is cited by key.
\begin{remark}[relation to the causal theory]\label{rem:provenance-causal}
(A1), (A2), (A4)--(A8) and (ND) are the axioms of the time-causal article
\sscite{fagerstrom2026hemigroup} up to the change of letters, and (A3) replaces causality by
reflection symmetry: the axiom that removed the Gaussian there is the axiom that admits it
here. Three consequences are visible already. The convolution representation yields symmetric
probability measures on $\RR$ rather than probability measures on a half-line, and the Fourier
transform can vanish where the Laplace transform cannot, so a nonvanishing lemma
(\ref{lem:nonvanishing}) replaces the positivity of the causal integrand. The drift slot of the
causal profile representation is occupied here by the Gaussian coefficient, so the member that
stands where the pure delay stood is the Gaussian scale space itself, admissible, extreme and
nondegenerate. And the gauge argument, which used the causal vanishing lemma twice, is
re-routed through a smoothed transform (\ref{cor:smoothed-transmittance}) and a dilation
argument (\ref{lem:dilation-atom}), the exclusion of lattice laws following from covariance
(\ref{lem:no-lattice}) rather than from the function class.
\end{remark}

% PROVED (proof 2026-09-09): both clauses' axiom content is machine-checked and moved to
% SpatialLine --- two_members_gaussian (SpatialLine/Gaussian.lean) and two_members_matern
% (SpatialLine/Matern.lean). The proof below is rewritten to be the proof of record for what was
% actually verified, and two of its steps changed:
%  * (A7). The old text derived strong convergence of the operators from continuity of the
%    transform "by the argument written out in the proof of thm:main-characterization". That is
%    an epsilon/3 argument and it is not needed: for the Gaussian the semigroup makes a
%    difference of operators factor through one increment, and for the Matern hemigroup the two
%    endpoints are moved separately through an intermediate. Both then reduce to the estimate
%    || mu*f - f ||_1 <= int Theta_f d mu, which is Fubini. The node no longer depends on the
%    continuity argument inside thm:main-characterization, and clause (2)'s axioms no longer go
%    through that theorem's constructive direction at all.
%  * (ND). The old text exhibited a frequency at which the transform is not 1. The axiom is
%    about the OPERATOR, so what is needed is that convolution by a probability measure is the
%    identity only for delta_0 --- proved by testing at a bump with a strict maximum at the
%    origin. That lemma (SpatialLine.eq_dirac_of_mconvL1_eq_id) is what
%    lem:convolution-representation's uniqueness clause should start from.
% No ledger interface is consumed: in particular the Matern clause never reaches the Bessel-K
% density (A15/A16), and prop:levy-continuity's first clause (A6) is proved from Mathlib.
% The node stays \notready: two_members_matern_moments and matern_density are still sorry.
% SKELETON NOTE (phase B, 2026-09-09): clause (2)'s sentence "its kernel at scale t is the
% Matern kernel of smoothness gamma - 1/2 and range t" is discharged by Skeleton.matern_density
% (Chapter 10), which is now named in the \lean tag below. Phase A's question Q1 asked whether
% to add the declaration or to move the sentence into the annotation; the declaration is the
% more faithful tag and the forward \uses edges to prop:matern-density were already in place,
% so the statement text is unchanged.
% CHANGED (2026-09-14, R156; a NARROWING, requested by the article session on 2026-09-12 and
% required by the external review's finding C): the density sentence LEAVES clause (2). The
% density (|x|/2t)^{gamma-1/2} K_{gamma-1/2}(|x|/t) is a probability density for every gamma > 0,
% but a Matern COVARIANCE function only for smoothness gamma - 1/2 > 0: at gamma = 1/2 it is K_0,
% which diverges logarithmically at the origin, so the identification as stated was false on
% 0 < gamma <= 1/2, which this node's own hypothesis admits. Three consequences, all in the
% direction of claiming less: the sentence goes (the mean, the variance and the finiteness of
% every moment stay); Skeleton.matern_density leaves the \lean tag and prop:matern-density leaves
% \uses; and two_members_matern_moments takes its integrability conjunct from matern_moments
% (the even moments of the Gamma mixture with |x|^n <= 1 + x^{2n}) instead of
% matern_moments_integrable (the moment criterion, ledger A13). So the node is \leanok and prints
% Lean core, and the closed form of the density is a cited fact of chapter 10 that nothing here
% asserts. The paper's copy of the statement is updated verbatim; the draft is mirrored.
% draft: the two members of \S3 ("Two members, at the two corners of what the axioms admit")
\begin{proposition}[two members]\label{prop:two-members}
\uses{def:cascade-family,prop:matern-exponent,thm:main-characterization,lem:convolution-representation,prop:fourier-uniqueness,prop:levy-continuity}
\lean{SpatialLine.two_members_gaussian, SpatialLine.two_members_matern,
SpatialLine.two_members_matern_moments}\leanok
\notes{gamma-matern-corner}
\begin{enumerate}
\item \emph{(The Gaussian family.)} Let $g_v$ be the centred Gaussian law of variance $v$, so
that $\hat g_v(\omega) = e^{-v\omega^2/2}$, and put $\Phis{s}{t}f := g_{t^2-s^2} * f$ for
$0 \le s \le t$. This family satisfies (A1)--(A8) and (ND), with $S_\lambda t = \lambda t$.
\item \emph{(The \Matern{} family.)} For $\gamma > 0$ let $\rho_{s,t}$ be the symmetric
probability measure with
\[
  \hat\rho_{s,t}(\omega) = \Bigl(\frac{1 + s^2\omega^2}{1 + t^2\omega^2}\Bigr)^{\gamma},
  \qquad 0 \le s \le t,
\]
and put $\Phis{s}{t}f := \rho_{s,t} * f$. This family satisfies (A1)--(A8) and (ND), with
$S_\lambda t = \lambda t$. Its kernel at scale $t$ has mean displacement $0$, variance
$2\gamma t^2$, and every moment finite.
\end{enumerate}
\statusT\quad\emph{(The Gaussian clause is elementary at the kernel: variances add, which is
(A6); the image of $g_v$ under $x \mapsto \lambda x$ is $g_{\lambda^2v}$, which is (A8); and
$g_v$ being a symmetric probability density gives (A1)--(A5). The \Matern{} clause is elementary
in the same way once its kernels are known to be symmetric probability measures, which is the
one non-obvious point and is \ref{prop:matern-exponent}. The forward reference is deliberate:
this proposition is where the article says what the theory must fit between, while the existence
of the second member is a corner computation belonging to Chapter~\ref{ch:corners}. Two clauses
are not elementary at the kernel and are proved at the operator, as the proof shows: (A7), which
asks for strong convergence in $\Lone$ and not merely weak convergence of the kernels, and (ND),
which asks that the operator is not the identity and not merely that the transform is not $1$.
\textbf{Both members are machine-checked in full, and the node prints Lean core}: the moment
sentence of (2) is \texttt{two\_members\_matern\_moments}, whose mean and variance are
\ref{prop:moments-tails}(1)'s variance formula on this profile --- the catalogue's first moment
being $\int_0^\infty x\,2\gamma e^{-x}dx = 2\gamma$ --- and whose finiteness of every moment is
the even moments of the Gamma mixture of \ref{prop:matern-exponent}(3) with
$|x|^n \le 1 + x^{2n}$, so that the moment criterion of \ref{prop:moments-tails}, ledger A13, is
no longer on the path (2026-09-14, R156). What the node no longer says is the \emph{density}
sentence: the identification of the kernel at scale $t$ as the \Matern{} kernel of smoothness
$\gamma - \tfrac12$ and range $t$, which holds as a covariance function only for
$\gamma > \tfrac12$ and is \ref{prop:matern-density}, cited there and asserted here of no
member.)}
\end{proposition}

\begin{proof}
\leanok
(1) Each $g_v$ with $v > 0$ is a symmetric probability density, and $g_0 := \delta_0$; so
(A1)--(A5) hold by the converse clause of \ref{lem:convolution-representation}, (A3) because
$g_v$ is invariant under reflection. For (A6),
$\hat g_{t^2-s^2}\,\hat g_{s^2-r^2} = \hat g_{t^2-r^2}$ and transforms determine measures
(\ref{prop:fourier-uniqueness}); $\Phis{t}{t} = \mathrm{Id}$ is $g_0 = \delta_0$. For (A8) with
$S_\lambda t = \lambda t$, the image of $g_v$ under $x \mapsto \lambda x$ is $g_{\lambda^2 v}$,
and $\lambda^2(t^2 - s^2) = (\lambda t)^2 - (\lambda s)^2$. For (ND): a probability measure whose
convolution operator fixes every $f \in \Lone$ is $\delta_0$ --- convolve a bump attaining a
strict maximum at the origin, and read the resulting identity there --- and $g_{t^2-s^2} =
\delta_0$ only when $s = t$.

For (A7) the semigroup does the work. For $0 \le v \le w$ one has $g_w = g_v * g_{w-v}$, so
$g_w * f - g_v * f = g_v * (g_{w-v} * f - f)$, and convolution by a probability measure is a
contraction of $\Lone$:
\[
  \| g_w * f - g_v * f \|_1 \ \le\ \| g_{w-v} * f - f \|_1 .
\]
Continuity in $(s,t)$ therefore reduces to continuity of $v \mapsto g_v * f$ at $v = 0$, and that
follows from the elementary estimate
\[
  \| \mu * f - f \|_1 \ \le\ \int_\RR \Theta_f(y)\, \mu(dy),
  \qquad \Theta_f(y) := \| \Tr{y} f - f \|_1,
\]
valid for every probability measure $\mu$, which is Fubini applied to
$(\mu * f - f)(x) = \int (f(x-y) - f(x))\,\mu(dy)$. The function $\Theta_f$ is continuous,
bounded by $2\|f\|_1$, and vanishes at $y = 0$, by strong continuity of translation on $\Lone$;
and $g_v$ is the image of $g_1$ under $z \mapsto \sqrt{v}\,z$, so
$\int \Theta_f\,dg_v = \int \Theta_f(\sqrt{v}\,z)\,g_1(dz) \to 0$ as $v \to 0$ by dominated
convergence.

(2) By \ref{prop:matern-exponent} the function $F(\omega) = \gamma\log(1 + \omega^2)$ is
admissible in the sense of \ref{thm:main-characterization}, with Gaussian coefficient $0$ and
folded profile $k(x) = 2\gamma e^{-x}$, and
$\hat\rho_{s,t}(\omega) = \exp[-(F(t\omega) - F(s\omega))]$; so $\rho_{s,t}$ is the increment
kernel of the admissible family attached to $F$ in the canonical gauge. The mean, the variance
and the finiteness of every moment are \ref{prop:matern-exponent}(3): the mean and the variance
from the variance formula on the profile $2\gamma e^{-x}$, and every moment from the even
moments of the Gamma mixture, $|x|^n \le 1 + x^{2n}$ bounding the odd ones.

The axioms themselves are read off the transform and need neither of those. (A1)--(A5) are again
\ref{lem:convolution-representation}, the kernels being symmetric probability measures; (A6) is
$\hat\rho_{r,s}\,\hat\rho_{s,t} = \hat\rho_{r,t}$ together with $\hat\rho_{t,t} \equiv 1$ and
\ref{prop:fourier-uniqueness}; (A8) with $S_\lambda t = \lambda t$ is
$\hat\rho_{s,t}(\lambda\omega) = \hat\rho_{\lambda s,\,\lambda t}(\omega)$; and (ND) is
$\hat\rho_{s,t}(1) = \bigl((1+s^2)/(1+t^2)\bigr)^{\gamma} < 1$ for $s < t$.

(A7) costs more than in (1), because a hemigroup has no single increment parameter: $\Phis{s}{t}$
and $\Phi_{s',t'}$ are compared through the intermediate $\Phi_{s',t}$, one endpoint at a time.
Moving the left endpoint leaves an increment applied \emph{first}, which the contraction bound
absorbs as in (1); moving the right endpoint leaves one applied \emph{last}, to $\Phi_{s',t}f$
rather than to $f$, so the estimate above is read at that element --- which loses nothing, since
translation commutes with convolution and hence $\Theta_{\mu * f} \le \Theta_f$. Both remainders
are increments whose two endpoints come together, and for those $\hat\rho_{a,b} \to 1$
pointwise, so $\rho_{a,b} \to \delta_0$ weakly (\ref{prop:levy-continuity}) and
$\int \Theta_f\,d\rho_{a,b} \to \Theta_f(0) = 0$. On the diagonal $s = t$ the intermediate is
unavailable and unnecessary: the limit operator is the identity, and one collapsing increment
suffices.

Both families are verified in the machine-checked development, every axiom including (A7); the
forward reference to \ref{prop:matern-exponent} carries the \emph{existence} of the \Matern{}
kernels and their moments, and is not used in the verification of the axioms.
\end{proof}

The Gaussian family plays here the role the pure delay played in the causal theory: the member
with no catalogue, the boundary the class closes on, and an extreme ray of the cone
(\ref{prop:choquet-cone}). Unlike the pure delay it is not degenerate. Between it and the
\Matern{} family the theory has to fit, and \ref{thm:main-characterization} says exactly what
fits.

% draft: Proposition 3.6
% CHANGED (skeleton 2026-09-08): the example's q is now required to have MEAN ZERO. Without it
% phi_eps(omega) = eps*int q - eps*qhat(omega) tends to eps*int q at infinity, and by
% Riemann--Lebesgue no member of A(R) does that -- so the exhibited family was not an instance
% of the proposition's own setting. The witness is unaffected: q = 1_{[1,2] u [-2,-1]} -
% 1_{[3,4] u [-4,-3]} is even, carried by {|x| >= 1}, of mean zero, and not a.e. nonnegative.
% Found by writing the Lean statement (Skeleton.no_positivity_example).
% CHANGED (review 2026-09-09, R28): the proof's last paragraph now says why BOTH signs of eps
% are excluded. Mean zero does two things, not one: it puts phi_eps in the Wiener algebra, and
% together with "not a.e. nonnegative" it gives "not a.e. nonpositive", which is what makes
% eps*q a non-positive measure for eps < 0 as well. The statement is unchanged.
% CHECK (draft, checks for \S2--\S3, item 5): the one imported ingredient in the proof was
% uniqueness of the SIGNED \Levy--Khintchine representation, which prop:fourier-toolbox(3)
% does NOT supply -- that clause is stated for positive measures. RESOLVED (proof 2026-09-10):
% the proposition does not need it. See the CHANGED note below.
% CHANGED (proof of record 2026-09-10): the example's proof is rewritten to the machine-checked
% route (SpatialLine.no_positivity_example). The printed route proved signed uniqueness inline
% and then read the example off it; the checked route averages both readings of F over [0,T]
% and reads the resulting identity twice -- at rate T^2, which fixes the Gaussian coefficient
% at 1, and at rate 1, where Fatou kills the \Levy{} measure -- so no signed measure is ever
% transformed and prop:fourier-uniqueness is spent only on the Jordan parts of q, exactly as
% wave 3 had already proved. Two things the rewrite records: the finiteness of the \Levy{} measure,
% which the surveyed route established as a separate step, is not needed (Fatou asks nothing of
% it); and the statement is unchanged. The signed-uniqueness argument survives as
% rem:signed-levy-uniqueness, marked not machine-checked, because it is the more informative
% statement.
% CHANGED (proof of record 2026-09-10, wave 6): the POSITIVITY paragraph is rewritten to the
% machine-checked route (SpatialLine.no_positivity_classification), and both directions gain a
% step the printed argument did not have.
% Forwards, the printed proof says "the family satisfies every axiom, so
% thm:main-characterization applies". It does not yet: that theorem is stated for PROBABILITY
% kernels and this node's kernels are integrable FUNCTIONS not known to be nonnegative -- which
% is what the clause is about. The missing step is where (A4) is spent, and it is the
% approximate identity plus the closedness of the positive cone of L^1. The printed proof also
% says the gauge is the identity "since S_lam t = lam t". The checked route does not prove that
% and does not need it: it reads the conclusion at t = 1 and absorbs the constant chi(1) by
% DILATING the profile (SpatialLine.SDProfile.dilate), which is one definition against a
% rigidity argument. This is also the one place F(0) = 0 is spent, which is why the split
% auxiliary needed that normalisation and the node does not (it has phi(0) = 0 by hypothesis).
% Backwards, the printed proof says the profile form "makes every kernel a probability measure",
% which again presumes the identification of mu_{s,t} with the constructed law. The checked
% route makes it: the Jordan split against nu + mu^- dx, one application of
% prop:fourier-uniqueness, and one test of the resulting identity on the set where mu_{s,t} is
% negative. That is the signed-uniqueness step the CHECK above records as not supplied by
% prop:fourier-toolbox(3), and it is supplied here rather than imported
% (SpatialLine/SignedUniqueness.lean).
% CHANGED (2026-09-11, R122; a proof repair, no statement or status changed): the printed proof
% said b -> phi(b .) is continuous into A(R) and derived (A7) from the norm continuity of h_{s,t}.
% That is false at the edge s = 0: for b > 0, phi(b .) is the transform of the dilate h_b, whose
% L^1 norm is ||h||_1, while phi(0 .) = 0, so the map jumps at b = 0 unless phi = 0. (A7) holds
% all the same, by strong continuity of the convolution operators h_b * ., which is what the proof
% now argues. Found by the estimate of no_positivity_family, which priced a Lean route along the
% same operator form at about 615 lines and was not attempted (over the pass's 600-line bound).
% Draft Proposition 3.6's proof mirrored.
% CHANGED (article mirror 2026-09-14, R157; a proof repair, no statement or status changed): the
% Gaussian factor of the construction had the wrong variance. The proof needs a kernel with
% transform e^{-(t^2-s^2)omega^2}, and under this article's normalisation ghat_v = e^{-v omega^2/2}
% (prop:two-members) that is the Gaussian of variance 2(t^2-s^2), not of variance t^2-s^2. Found by
% the external review (finding A, 2026-09-14); it is the one construction of the article that is not
% machine-checked, so no declaration carried the factor. The normalisation is now named where the
% kernel is introduced.
% CHANGED (article mirror 2026-09-14, R157): retitled. "Without positivity there is no
% classification" reads as an impossibility theorem about every possible classification, which is
% more than the node proves and more than the article claims; the external review read it that way.
% The title now says what the statement exhibits, and the impossibility is where it belongs, in the
% annotation and the lead, stated of a classification in the style of thm:main-characterization.
% The label is unchanged. The paper's subsection heading carries the same wording.
\begin{proposition}[without positivity: an affine family through the Gaussian]
\label{prop:no-positivity-no-classification}
\uses{def:cascade-family,thm:main-characterization,lem:selfdecomposable-exponents,prop:fourier-uniqueness,prop:fourier-toolbox}
\lean{Skeleton.no_positivity_no_classification, Skeleton.no_positivity_family,
SpatialLine.no_positivity_classification, SpatialLine.no_positivity_example}\notready
Let $A(\RR) := \{\hat h : h \in \Lone\}$ be the Wiener algebra and let $\phi \in A(\RR)$ be
real and even with $\phi(0) = 0$. Put $F(\omega) := \omega^2 + \phi(\omega)$ and define
$\mu_{s,t}$, $0 \le s \le t$, by $\hat\mu_{s,t}(\omega) = e^{-(F(t\omega) - F(s\omega))}$. Then
each $\mu_{s,t}$ is a finite symmetric signed measure, an integrable function when $s < t$, and
$\Phis{s}{t} := \mu_{s,t} * {}$ satisfies (A1)--(A3), (A5)--(A8) with
$S_\lambda t = \lambda t$, and (ND). It satisfies (A4) if and only if $F$ is of the form
\ref{eq:sd-profile}. In particular, for $q \in \Lone$ even, supported in $\{|x| \ge 1\}$, of
mean zero and not almost everywhere nonnegative, and
$\phi_\varepsilon(\omega) := \varepsilon\int_\RR (1 - \cos\omega x)\,q(x)\,dx$, the family
violates (A4) for every $\varepsilon \ne 0$.
\statusT\quad\emph{(The inversion of the semigroup literature's headline, and the article's
thesis in negative form: under the hemigroup the non-positive class contains an
infinite-dimensional affine family through the Gaussian, so no classification in the style of
\ref{thm:main-characterization} is possible without (A4). No attempt is made to describe the
non-positive class beyond this. Every step is elementary. \textbf{Partly formalised}
(2026-09-10): the final sentence, that the family of the example violates (A4) for every
$\varepsilon \ne 0$, is machine-checked as \texttt{SpatialLine.no\_positivity\_example} and
rests on Lean core; so does the Ces\`aro identity it is built from, which is stated for an
arbitrary symmetric \Levy{} pair. \textbf{The positivity equivalence is machine-checked as well}
(2026-09-10, wave 6), as \texttt{SpatialLine.no\_positivity\_classification}, by the two routes
this annotation had predicted: forwards by reading (A4) at an approximate identity, which makes
the function kernel the density of a probability measure and hands the family to
\ref{thm:main-characterization}; backwards by identifying $k\,dx$ with the constructed law
through the Jordan split. It spends \ledger{A1} and \ledger{A3}, and each direction spends less
than the pair: the forward one only \ledger{A3}'s uniqueness clause, the backward one only
\ledger{A1} with \ledger{A3}'s converse. \textbf{One clause is unchecked}, the Wiener-algebra
construction of $h_{s,t}$: Mathlib carries the $\Lone$ Fourier transform but not $A(\RR)$ as a
Banach algebra, and that is the whole of what keeps this node \texttt{notready}. \textbf{Priced
2026-09-11}: about $615$ Lean lines (range $480$--$720$), on Lean core, by a route that needs no
Wiener algebra --- exponentiate the convolution operators $f \mapsto h_b * f$ in the Banach
algebra of bounded operators on $\Lone$ rather than the functions, and read every transform
through the character pairing, as the proof below does for (A7). Not attempted in that pass,
being over its bound; the pieces, the risks and a stop-loss are recorded with the estimate.)}
\end{proposition}

\begin{proof}
\emph{The family.} $A(\RR)$ is a Banach algebra under pointwise multiplication, invariant
under dilation, so $\phi_{s,t} := \phi(t\,\cdot) - \phi(s\,\cdot)$ lies in it, and
$e^{-\phi_{s,t}} = 1 + \hat h_{s,t}$ with $h_{s,t} \in \Lone$, the exponential series converging
in the algebra. (The map $b \mapsto \phi(b\,\cdot)$ is \emph{not} norm-continuous into $A(\RR)$
at $b = 0$: for $b > 0$ it is the transform of the dilate $h_b := b^{-1}h(\cdot/b)$, of norm
$\|h\|_1$, while $\phi(0\,\cdot) = \phi(0) = 0$. So $h_{s,t}$ is not norm-continuous up to the
edge $s = 0$, and (A7) is argued below through operators.) With $g_{2(t^2-s^2)}$ the
Gaussian kernel of variance $2(t^2-s^2)$, whose transform is $e^{-(t^2-s^2)\omega^2}$ in the
normalization $\hat g_v = e^{-v\omega^2/2}$ of \ref{prop:two-members} --- a density for
$s < t$, and $\delta_0$ on the diagonal
--- we get $\mu_{s,t} = g_{2(t^2-s^2)} + g_{2(t^2-s^2)} * h_{s,t}$: a finite signed measure,
symmetric because its transform is real and even, and an integrable function for $s < t$.

\emph{The axioms.} (A1) holds with $\|\mu_{s,t}\| \le 1 + \|h_{s,t}\|_1$; (A2) because
convolution commutes with translation; (A3) by symmetry; (A5) from $\hat\mu_{s,t}(0) = 1$; (A6)
because the exponents add and transforms determine measures (\ref{prop:fourier-uniqueness});
(A8) with $S_\lambda t = \lambda t$, because $F(\lambda t\omega) - F(\lambda s\omega)$ is
$F(t\,\cdot) - F(s\,\cdot)$ evaluated at $\lambda\omega$; and (ND) because
$F(t\omega) - F(s\omega) = (t^2-s^2)\omega^2 + \phi_{s,t}(\omega)$ is unbounded, $\phi_{s,t}$
being bounded. For (A7), let $h_e$ be the even part of $h$ and $H_b f := h_b * f$ for
$b > 0$, with $h_b$ the dilate of $h_e$, and $H_0 := 0$. Each $H_b$ is a bounded operator on
$\Lone$ of norm at most $\|h\|_1$, the $H_b$ commute with one another and with the Gaussian
convolutions, and $b \mapsto H_b f$ is continuous on $[0,\infty)$ for every $f \in \Lone$: at
$b > 0$ by strong continuity of dilation, and at $b = 0$ because $h_b$ is an approximate
identity of total mass $\int h_e = \phi(0) = 0$, so $h_b * f \to 0$ in $\Lone$. Then
$\mu_{s,t} * f = g_{2(t^2-s^2)} * \exp\bigl(-(H_t - H_s)\bigr)f$, the exponential series converging in
operator norm uniformly in $(s,t)$, its $n$-th term being bounded by $(2\|h\|_1)^n/n!$. A product
of uniformly bounded, strongly continuous operator families is strongly continuous, so each
partial sum is, hence so is the sum, and composing with the strongly continuous Gaussian family
gives (A7). The route is strong continuity throughout, and it does not need $h_{s,t}$ to depend
continuously on $(s,t)$ in norm, which it does not at $s = 0$.

\emph{Positivity.} Suppose (A4) holds. The kernels $\mu_{s,t}$ are so far only integrable
functions, and \ref{thm:main-characterization} is a statement about probability kernels, so the
first step is to see that (A4) makes them nonnegative. Convolve with an approximate identity:
$\mu_{s,t} * \rho_\varepsilon \ge 0$ for every $\varepsilon > 0$ because $\rho_\varepsilon \ge 0$,
and $\mu_{s,t} * \rho_\varepsilon \to \mu_{s,t}$ in $\Lone$; the positive cone of $\Lone$ is
closed, so $\mu_{s,t} \ge 0$ almost everywhere. Its total mass is $\hat\mu_{s,t}(0) = 1$, so
$\mu_{s,t}\,dx$ is a probability measure, the family satisfies every axiom, and
\ref{thm:main-characterization} applies. It returns a gauge $\chi$ and a profile $P$ with
$F(t\omega) - F(0) = P(\chi(t)\omega)$; reading that at $t = 1$ and using $F(0) = \phi(0) = 0$
gives $F(\omega) = P(\chi(1)\omega)$, and $\chi(1) > 0$, so $F$ is $P$ dilated and is itself of
the form \ref{eq:sd-profile}.

Conversely, suppose $F$ has that form. Each increment exponent
$\omega \mapsto F(t\omega) - F(s\omega)$ is then a member of $\NDs$ by
\ref{lem:selfdecomposable-exponents}, hence the exponent of a symmetric probability law $\nu$
(\ref{prop:fourier-toolbox}); and $\nu$ has the same cosine transform as $\mu_{s,t}$, which is
even and integrable. The Jordan parts of $\mu_{s,t}$ split it as a difference of two finite
measures whose characteristic functions are determined by that cosine transform, evenness
making both sine transforms vanish, so $\mu_{s,t}^+\,dx = \nu + \mu_{s,t}^-\,dx$ by
\ref{prop:fourier-uniqueness}. Testing this on the set where $\mu_{s,t}$ is negative makes its
left side vanish, so $\mu_{s,t}^- = 0$ almost everywhere: $\mu_{s,t} \ge 0$, and (A4) is
positivity of an integral.

\emph{The example.} Write $g(\omega) := \int (1 - \cos\omega x)\,q(x)\,dx$, so that
$F = \omega^2 + \varepsilon g$. Since $\int q = 0$ we have $g = -\hat q$, so $g$ is bounded and
$g(0) = 0$; hence $\phi_\varepsilon = -\varepsilon\,\hat q$ lies in $A(\RR)$, real and even
because $q$ is, and the family is one of those described. (Mean zero is not cosmetic: without
it $\phi_\varepsilon$ tends to $\varepsilon\int q \ne 0$ at infinity, which no member of the
Wiener algebra does.)

Suppose $F$ were of the form \ref{eq:sd-profile}, with Gaussian coefficient $a \ge 0$ and
\Levy{} measure $\nu$ on $(0,\infty)$ subject to $\int (1 \wedge x^2)\,\nu(dx) < \infty$.
Average the two readings of $F$ over $[0,T]$. The elementary primitive
$\int_0^T (1 - \cos\omega x)\,d\omega = T\bigl(1 - \frac{\sin Tx}{Tx}\bigr)$ and Tonelli give,
for every $T > 0$,
\[
  \frac{T^2}{3}
   + \varepsilon\int \Bigl(1 - \frac{\sin Tx}{Tx}\Bigr) q(x)\,dx
  \;=\; \frac{a\,T^2}{3}
   + \int \Bigl(1 - \frac{\sin Tx}{Tx}\Bigr) \nu(dx),
\]
both integrals converging absolutely because
$0 \le 1 - \frac{\sin u}{u} \le \min(2,\, u^2/6)$.

Read this identity at rate $T^2$. Divided by $T^2$, the perturbation term vanishes in the limit
because its integral is bounded by $2\|q\|_1$, and the \Levy{} term vanishes by dominated
convergence, the integrand $T^{-2}\bigl(1 - \frac{\sin Tx}{Tx}\bigr)$ being at most
$2\,(1 \wedge x^2)$ for $T \ge 1$ and tending to $0$ pointwise. Hence $a = 1$, and the identity
reduces to
$\int \bigl(1 - \frac{\sin Tx}{Tx}\bigr)\nu(dx)
 = \varepsilon\int \bigl(1 - \frac{\sin Tx}{Tx}\bigr) q(x)\,dx$.

Now read it at rate $1$. The right-hand side tends to $\varepsilon \int q = 0$ as
$T \to \infty$, by dominated convergence against $|q|$. On the left,
$1 - \frac{\sin Tx}{Tx} \to 1$ at every $x \ne 0$ and $\nu$ gives the origin no mass, so Fatou
forces $\nu\bigl((0,\infty)\bigr) = 0$. With $a = 1$ and $\nu = 0$, \ref{eq:sd-profile} reads
$F(\omega) = \omega^2$, that is $\varepsilon\,g \equiv 0$, that is $\hat q \equiv 0$ because
$\varepsilon \ne 0$. An even integrable function whose cosine transform vanishes identically is
null: its positive and negative parts are two finite measures whose cosine transforms agree and
whose sine transforms both vanish, evenness being what kills the latter, so
\ref{prop:fourier-uniqueness} identifies them. Thus $q = 0$ almost everywhere, against the
hypothesis that $q$ is not almost everywhere nonnegative --- and both signs of $\varepsilon$
are covered, neither the identity nor either reading of it seeing the sign. So $F$ is not of
the form \ref{eq:sd-profile}, and (A4) fails by the previous paragraph.
\end{proof}

\begin{remark}[the route through signed uniqueness]\label{rem:signed-levy-uniqueness}
The last clause of the proposition, that the family built from $q$ violates (A4), is usually
argued the other way round: $\omega^2 + \phi_\varepsilon$ \emph{is} a
signed \Levy--Khintchine representation, with Gaussian coefficient $1$ and signed measure
$\varepsilon\,q(x)\,dx$, and such representations are unique among signed measures $m$ with
$\int (1 \wedge x^2)\,|m|(dx) < \infty$. Since $\varepsilon\,q$ is not a nonnegative measure,
$F$ is not of the form \ref{eq:sd-profile}. The uniqueness is the averaging identity
\[
  \frac{1}{2h}\int_{\omega-h}^{\omega+h}\psi(u)\,du \;-\; \psi(\omega)
  \;=\; \frac{a\,h^2}{3}
   + \int \cos(\omega x)\Bigl(1 - \frac{\sin hx}{hx}\Bigr)\,m(dx)
\]
for $\psi(\omega) = a\omega^2 + \int (1-\cos\omega x)\,m(dx)$, whose right-hand side is a
constant plus the Fourier--Stieltjes transform of the finite signed measure
$\bigl(1 - \frac{\sin hx}{hx}\bigr)\,m(dx)$; so $\psi$ determines that measure for every
$h > 0$, hence $m$ off the origin, hence $a$. This is the more informative statement, because
it is about the representation rather than about one $F$. It is
also strictly more than the proposition needs, and it is \textbf{not machine-checked}: the
proof of record above replaces it by the two readings of the averaged identity, at the rates
$T^2$ and $1$, which identify $a$ and $\nu$
for the single exponent at hand and never mention a signed measure.
\end{remark}

Under recursivity the same setting forces $F(\omega) = a|\omega|^\alpha$ with $\alpha > 0$: a
one-parameter class whether or not positivity is imposed, positivity entering only to bound
$\alpha \le 2$ (\sscite{pauwels1995extended}, Prop.~1). Under the hemigroup the non-positive
class contains the affine family $\omega^2 + A(\RR)_{\mathrm{even},0}$ through the Gaussian ---
and the same family through every admissible exponent unbounded at infinity --- so it is
positivity, not recursivity, that narrows.
